% figures/fig4_geometria_to_a.tex
\begin{tikzpicture}[
    >=Stealth,
    stand/.style={
        line width=1.5pt, draw=black!75
    },
    baseplate/.style={
        draw=black!80, fill=black!25, rounded corners=1pt
    },
    sensor/.style={
        circle, draw=black!85, fill=black!12, line width=1.1pt,
        minimum size=0.60cm, inner sep=0pt
    },
    % Fuente en naranja unificado
    source/.style={
        circle, draw=black!85, fill=orange!80, line width=1.1pt,
        minimum size=0.60cm, inner sep=0pt
    },
    tabletop/.style={
        draw=black!75, line width=0.9pt, pattern=north east lines, pattern color=black!25
    },
    raydir/.style={
        ->, line width=1.2pt, draw=UdeAGreen!75!black
    },
    rayrefl/.style={
        ->, line width=1.1pt, dashed, draw=red!70!black
    }
]

    % =========================================================================
    % 1. SUPERFICIE HORIZONTAL DE LA MESA (y = 0)
    % =========================================================================
    \draw[tabletop] (-5.6, -0.38) rectangle (5.6, 0.0);
    \node[below=4pt, font=\footnotesize\sffamily] at (0, -0.38) {Superficie reflectora horizontal (Mesa / Suelo)};

    % =========================================================================
    % 2. SOPORTES Y TRANSDUCTORES (x = -3.6 y x = +3.6, y = 2.6)
    % =========================================================================
    \coordinate (Mic) at (-3.6, 2.6);       % Receptor R
    \coordinate (Spk) at (3.6, 2.6);        % Emisor S
    \coordinate (Bounce) at (0.0, 0.0);     % Punto de rebote en la mesa

    % Soporte y cápsula Receptor R (Gris)
    \draw[baseplate] (-4.05, 0.0) rectangle (-3.15, 0.14);
    \draw[stand] (-3.6, 0.14) -- (Mic);
    \node[sensor] (NodeMic) at (Mic) {};
    \node[left=8pt, font=\small\sffamily\bfseries] at (Mic) {Receptor ($R$)};

    % Soporte y cápsula Fuente S (Naranja)
    \draw[baseplate] (3.15, 0.0) rectangle (4.05, 0.14);
    \draw[stand] (3.6, 0.14) -- (Spk);
    \node[source] (NodeSpk) at (Spk) {};
    \node[right=8pt, font=\small\sffamily\bfseries] at (Spk) {Fuente ($S$)};

    % Arcos acústicos de radiación en naranja unificado
    \draw[line width=0.85pt, orange!85!black] ($(Spk)+(-0.45, 0.3)$) arc (145:215:0.5);
    \draw[line width=0.85pt, orange!85!black] ($(Spk)+(-0.75, 0.45)$) arc (145:215:0.8);
    \draw[line width=0.85pt, orange!85!black] ($(Spk)+(-1.05, 0.6)$) arc (145:215:1.1);

    % =========================================================================
    % 3. TRAYECTORIAS ONDULATORIAS Y COTA CENTRAL DE ALTURA h
    % =========================================================================
    % Trayectoria directa horizontal
    \draw[raydir] (NodeSpk.west) -- node[above=4pt, font=\small\sffamily] {Onda directa: $r_{\text{dir}} = r$} (NodeMic.east);

    % Trayectoria reflejada especular en V
    \draw[rayrefl] (NodeSpk.south) -- node[below right=2pt, font=\footnotesize\sffamily, text=red!75!black] {$r_{\text{refl}}/2$} (Bounce);
    \draw[rayrefl] (Bounce) -- node[below left=2pt, font=\footnotesize\sffamily, text=red!75!black] {$r_{\text{refl}}/2$} (NodeMic.south);
    \fill[red!70!black] (Bounce) circle (2.0pt);

    % Eje vertical de altura h
    \draw[dashed, draw=black!45, line width=0.6pt] (0.0, 0.0) -- (0.0, 2.6);
    \draw[<->, draw=black!75, line width=0.65pt] (-0.35, 0.05) -- node[left=1pt, font=\small\sffamily] {$h$} (-0.35, 2.55);

    % =========================================================================
    % 4. CAJA METROLÓGICA (y = 5.0)
    % =========================================================================
    \node[draw=black!75, fill=white, rounded corners=4pt, line width=0.8pt, inner sep=8pt] at (0.0, 5.0) {
        \footnotesize
        \begin{tabular}{@{} l l @{}}
            \multicolumn{2}{@{} l @{}}{\textbf{Geometría de Reflexión y Condición Cuasi-Anecoica:}} \\[4pt]
            Trayectoria directa: & $r_{\text{dir}} = r$ \\
            Trayectoria reflejada: & $r_{\text{refl}} = 2\sqrt{(r/2)^2 + h^2} = \sqrt{r^2 + 4h^2}$ \\
            Diferencial de camino: & $\Delta r_{\text{refl}} = \sqrt{r^2 + 4h^2} - r$ \\[4pt]
            \multicolumn{2}{@{} l @{}}{\textcolor{UdeAGreen!80!black}{\textbf{Condición de corte:}} \ $c(T) \cdot T_{\text{gate}} < \Delta r_{\text{refl}} \implies \text{Anulación de eco }(c_{\text{room}} \to 0)$}
        \end{tabular}
    };

\end{tikzpicture}